% ===========================
% capitulo_implementacion.tex
% ===========================

\section{Módulo de Simulación}% chktex 24
\label{cap:bitacora-implementacion}

\subsection{Introducción}
Este capítulo documenta de forma técnica y reproducible el proceso de implementación del \emph{Módulo de Simulación} empleado para estudiar dinámicas gravitacionales de dos y tres cuerpos, así como la estimación del exponente de Lyapunov (LE) a partir de trayectorias numéricas. Se describe la arquitectura, el \emph{stack} tecnológico, y se ofrece una documentación por script (clases, funciones, parámetros, valores de retorno y ejemplos de uso).

\subsection{Arquitectura de Alto Nivel}
El módulo se divide en cuatro componentes principales:
\begin{enumerate}
    \item \textbf{Adaptador de dinámica gravitacional (\texttt{rebound\_adapter.py})}: encapsula la creación y propagación de simulaciones usando REBOUND, con saneamiento de entradas y extracción discreta de trayectorias.
    \item \textbf{Dinámica determinista (\texttt{dynamics.py})}: ofrece el campo vectorial (RHS) y el Jacobiano del sistema de dos cuerpos en 2D (CPU), y el esqueleto para aceleración GPU.\
    \item \textbf{Estimación de Lyapunov (\texttt{lyapunov.py})}: calcula el exponente máximo (mLCE) a partir de una simulación REBOUND (variacionales) o mediante un esquema discreto tipo Benettin con re-normalización periódica.
    \item \textbf{Inicialización del paquete (\texttt{\_\_init\_\_.py})}: expone una API mínima y controlada para importar las clases principales.
\end{enumerate}

\subsection{Stack tecnológico}
\begin{itemize}
    \item \textbf{Python 3.x} como lenguaje principal.
    \item \textbf{NumPy} para cálculo numérico denso.
    \item \textbf{REBOUND} para integración N-cuerpos (integradores \texttt{WHFast}, \texttt{IAS15}, etc.).
    \item \textbf{(Opcional) CuPy} para acelerar kernels en GPU (interfaz en \texttt{DynamicsGPU}).
\end{itemize}

% =========================================
% DOCUMENTACIÓN POR SCRIPT
% =========================================

% ---------- rebound_adapter.py ----------
\subsection{\texttt{rebound\_adapter.py} — Adaptador REBOUND}
\paragraph{Descripción general}
\texttt{ReboundSim} provee una interfaz ligera para: (i) crear una simulación a partir de masas, posiciones y velocidades iniciales; (ii) configurar integrador y sistema de unidades; (iii) integrar la dinámica y extraer una trayectoria discreta en forma de arreglo NumPy \([{\rm pasos}, {\rm cuerpos}, 6]\) con \((x,y,z,v_x,v_y,v_z)\).

\subsubsection{\texttt{class ReboundSim}}
\paragraph{Propósito} Gestionar la vida de una simulación REBOUND (creación, \emph{setup}, integración discreta).

\subsubsection*{\texttt{\_\_init\_\_(G:\ float = 1.0, integrator: str = ``whfast'')}}
\textbf{Propósito.} Inicializa constantes globales de la simulación (gravitación \(G\)) y el integrador numérico por defecto.

\begin{table}[H]
  \centering
  \caption{Parámetros de \texttt{ReboundSim.\_\_init\_\_}}
  \small
  \begin{tabularx}{\textwidth}{@{}p{0.27\textwidth}p{0.33\textwidth}X@{}}
    \toprule
    \textbf{Parámetro} & \textbf{Tipo} & \textbf{Descripción} \\
    \midrule
    \texttt{G} & \texttt{float} & Constante gravitacional (unidades consistentes con REBOUND). \\
    \texttt{integrator} & \texttt{str} & Nombre del integrador (\texttt{whfast}, \texttt{ias15}, etc.). \\
    \bottomrule
  \end{tabularx}
\end{table}

\begin{table}[H]
  \centering
  \caption{Valor de retorno de \texttt{ReboundSim.\_\_init\_\_}}
  \small
  \begin{tabularx}{\textwidth}{@{}p{0.35\textwidth}X@{}}
    \toprule
    \textbf{Tipo} & \textbf{Descripción} \\
    \midrule
    \texttt{None} & Constructor, no retorna valor. \\
    \bottomrule
  \end{tabularx}
\end{table}

\begin{listing}[H]
\caption{Bloque de Código de \texttt{ReboundSim}}
\begin{minted}[fontsize=\small, breaklines, breakanywhere]{python}
class ReboundSim:
    def __init__(self, G: float = 1.0, integrator: str = "whfast") -> None:
        self.G = G
        self.integrator = integrator
\end{minted}
\end{listing}

\begin{listing}[H]
\caption{Ejemplo de inicialización de \texttt{ReboundSim}}
\begin{minted}[fontsize=\small, breaklines, breakanywhere]{python}
from simulacion.rebound_adapter import ReboundSim
sim_api = ReboundSim(G=1.0, integrator="whfast")
\end{minted}
\end{listing}

\subsubsection*{\texttt{setup\_simulation~(masses, r0, v0, move\_to\_com=True) -> Any}}
\textbf{Propósito.} Construir e inicializar \texttt{rebound.Simulation} con masas y estado inicial (posiciones y velocidades), con opción de mover el sistema al centro de masa.

\begin{table}[H]
  \centering
  \caption{Parámetros de \texttt{setup\_simulation}}
  \small
  \begin{tabularx}{\textwidth}{@{}p{0.27\textwidth}p{0.33\textwidth}X@{}}
    \toprule
    \textbf{Parámetro} & \textbf{Tipo} & \textbf{Descripción} \\
    \midrule
    \texttt{masses} & \texttt{Iterable[float]} & Masas por cuerpo (\(\ge 2\) y \(\le 3\) en la configuración base). \\
    \texttt{r0} & \texttt{Iterable[Iterable[float]]} & Posiciones \((x,y,z)\) por cuerpo. \\
    \texttt{v0} & \texttt{Iterable[Iterable[float]]} & Velocidades \((v_x,v_y,v_z)\) por cuerpo. \\
    \texttt{move\_to\_com} & \texttt{bool} & Si \texttt{True}, traslada al centro de masa. \\
    \bottomrule
  \end{tabularx}
\end{table}

\begin{table}[H]
  \centering
  \caption{Valor de retorno de \texttt{setup\_simulation}}
  \small
  \begin{tabularx}{\textwidth}{@{}p{0.35\textwidth}X@{}}
    \toprule
    \textbf{Tipo} & \textbf{Descripción} \\
    \midrule
    \texttt{rebound.Simulation} & Simulación con partículas agregadas y lista para integrar. \\
    \bottomrule
  \end{tabularx}
\end{table}

\begin{listing}[H]
\caption{Bloque de Código de \texttt{setup\_simulation}}
\begin{minted}[fontsize=\small, breaklines, breakanywhere]{python}
    def setup_simulation(
        self,
        masses: Iterable[float],
        r0: Iterable[Iterable[float]],
        v0: Iterable[Iterable[float]],
        move_to_com: bool = True,
    ) -> Any:
        try:
            import rebound  # type: ignore
        except Exception as e:
            raise ImportError(
                "rebound is required for ReboundSim. Use --install-hint for help."
            ) from e

        mass_tuple = tuple(float(m) for m in masses)
        n_bodies = len(mass_tuple)
        if n_bodies < 2 or n_bodies > 3:
            raise ValueError("ReboundSim soporta entre 2 y 3 cuerpos en la configuracion base.")
\end{minted}
\end{listing}

\begin{listing}[H]
\caption{Ejemplo de \texttt{setup\_simulation}}
\begin{minted}[fontsize=\small, breaklines, breakanywhere]{python}
masses = (1.0, 1.0)
r0 = ((-1.0, 0.0, 0.0), (1.0, 0.0, 0.0))
v0 = ((0.0, -0.5, 0.0), (0.0, 0.5, 0.0))
sim = sim_api.setup_simulation(masses, r0, v0, move_to_com=True)
\end{minted}
\end{listing}

\subsubsection*{\texttt{integrate~(sim, t\_end: float, dt: float) -> Any}}
\textbf{Propósito.} Avanzar la simulación y extraer una trayectoria discreta de estados \((x,y,z,v_x,v_y,v_z)\) por cuerpo y paso temporal.

\begin{table}[H]
  \centering
  \caption{Parámetros de \texttt{integrate}}
  \small
  \begin{tabularx}{\textwidth}{@{}p{0.27\textwidth}p{0.33\textwidth}X@{}}
    \toprule
    \textbf{Parámetro} & \textbf{Tipo} & \textbf{Descripción} \\
    \midrule
    \texttt{sim} & \texttt{rebound.Simulation} & Contexto REBOUND válido. \\
    \texttt{t\_end} & \texttt{float} & Tiempo total a integrar. \\
    \texttt{dt} & \texttt{float} & Paso de integración (discretización de salida). \\
    \bottomrule
  \end{tabularx}
\end{table}

\begin{table}[H]
  \centering
  \caption{Valor de retorno de \texttt{integrate}}
  \small
  \begin{tabularx}{\textwidth}{@{}p{0.35\textwidth}X@{}}
    \toprule
    \textbf{Tipo} & \textbf{Descripción} \\
    \midrule
    \texttt{numpy.ndarray} & Arreglo \([\text{pasos}, \text{cuerpos}, 6]\) con posiciones y velocidades. \\
    \bottomrule
  \end{tabularx}
\end{table}

\begin{listing}[H]
\caption{Bloque de Código de \texttt{integrate}}
\begin{minted}[fontsize=\small, breaklines, breakanywhere]{python}
    def integrate(self, sim: Any, t_end: float, dt: float) -> Any:
        """Integra el sistema y devuelve la trayectoria como arreglo numpy [pasos, cuerpos, 6]."""
        import numpy as _np

        if not hasattr(sim, "particles"):
            raise ValueError("Se esperaba una instancia de rebound.Simulation.")

        sim.dt = dt
        steps = max(1, int(_np.floor(t_end / dt)))
        n_bodies = len(sim.particles)
        traj = _np.zeros((steps, n_bodies, 6), dtype=float)
        base_t = sim.t if hasattr(sim, "t") else 0.0
        for i in range(steps):
            t_target = base_t + (i + 1) * dt
            sim.integrate(t_target)
            for j, p in enumerate(sim.particles):
                traj[i, j, 0] = float(p.x)
                traj[i, j, 1] = float(p.y)
                traj[i, j, 2] = float(p.z)
                traj[i, j, 3] = float(p.vx)
                traj[i, j, 4] = float(p.vy)
                traj[i, j, 5] = float(p.vz)
        return traj
\end{minted}
\end{listing}

\begin{listing}[H]
\caption{Ejemplo de integración y extracción de trayectoria}
\begin{minted}[fontsize=\small, breaklines, breakanywhere]{python}
traj = sim_api.integrate(sim, t_end=10.0, dt=0.5)  # shape: [steps, bodies, 6]
print(traj.shape)
\end{minted}
\end{listing}

% ---------- dynamics.py ----------
\subsection{\texttt{dynamics.py} — Dinámica (CPU/GPU)}
\paragraph{Descripción general}
Define la dinámica del problema de dos cuerpos en 2D (\texttt{DynamicsCPU}) y un esqueleto para aceleración en GPU (\texttt{DynamicsGPU}). El estado se ordena como
\[
x = (x_1, y_1, v_{x1}, v_{y1},\; x_2, y_2, v_{x2}, v_{y2}).
\]

\subsubsection{\texttt{class DynamicsCPU}}
\paragraph{Propósito} Proveer el RHS \( \dot{x} = f(x,t) \) y el Jacobiano \(J=\partial f/\partial x\) en 2D.

\subsubsection*{\texttt{f~(x, t, m1, m2, G) -> Any}}
\textbf{Propósito.} Evaluar el campo vectorial (posiciones y aceleraciones Newtonianas suavizadas por \(\varepsilon\) numérico) para dos cuerpos.

\begin{table}[H]
  \centering
  \caption{Parámetros de \texttt{DynamicsCPU.f}}
  \small
  \begin{tabularx}{\textwidth}{@{}p{0.27\textwidth}p{0.33\textwidth}X@{}}
    \toprule
    \textbf{Parámetro} & \textbf{Tipo} & \textbf{Descripción} \\
    \midrule
    \texttt{x} & \texttt{array-like (8,)} & Estado \((x_1,y_1,v_{x1},v_{y1},x_2,y_2,v_{x2},v_{y2})\). \\
    \texttt{t} & \texttt{float} & Tiempo de evaluación (no autónomo opcional). \\
    \texttt{m1, m2} & \texttt{float} & Masas de los cuerpos. \\
    \texttt{G} & \texttt{float} & Constante gravitacional. \\
    \bottomrule
  \end{tabularx}
\end{table}

\begin{table}[H]
  \centering
  \caption{Valor de retorno de \texttt{DynamicsCPU.f}}
  \small
  \begin{tabularx}{\textwidth}{@{}p{0.35\textwidth}X@{}}
    \toprule
    \textbf{Tipo} & \textbf{Descripción} \\
    \midrule
    \texttt{numpy.ndarray} & Vector de dimensión 8 con \((\dot{x}_1,\dot{y}_1,\dot{v}_{x1},\dot{v}_{y1},\dot{x}_2,\dot{y}_2,\dot{v}_{x2},\dot{v}_{y2})\). \\
    \bottomrule
  \end{tabularx}
\end{table}

\begin{listing}[H]
\caption{Bloque de Código de \texttt{DynamicsCPU.f}}
\begin{minted}[fontsize=\small, breaklines, breakanywhere]{python}
class DynamicsCPU:
    """Calcula el RHS y jacobiano del sistema de dos cuerpos en 2D."""

    def f(self, x: Any, t: float, m1: float, m2: float, G: float) -> Any:
        import numpy as _np

        eps = 1e-12
        x1, y1, vx1, vy1, x2, y2, vx2, vy2 = x
        dx = x2 - x1
        dy = y2 - y1
        r2 = dx * dx + dy * dy + eps
        r = _np.sqrt(r2)
        r3 = r2 * r

        ax1 = G * m2 * dx / r3
        ay1 = G * m2 * dy / r3
        ax2 = -G * m1 * dx / r3
        ay2 = -G * m1 * dy / r3

        return _np.array([vx1, vy1, ax1, ay1, vx2, vy2, ax2, ay2], dtype=float)
\end{minted}
\end{listing}

\begin{listing}[H]
\caption{Ejemplo de uso de \texttt{DynamicsCPU.f}}
\begin{minted}[fontsize=\small, breaklines, breakanywhere]{python}
import numpy as np
from simulacion.dynamics import DynamicsCPU

dyn = DynamicsCPU()
state = np.array([-1.0, 0.0, 0.0, -0.5, 1.0, 0.0, 0.0, 0.5], dtype=float)
rhs = dyn.f(state, t=0.0, m1=1.0, m2=1.0, G=1.0)
\end{minted}
\end{listing}

\subsubsection*{\texttt{jac~(x, t, m1, m2, G) -> Any}}
\textbf{Propósito.} Calcular la matriz Jacobiana \(8\times8\) del sistema 2D de dos cuerpos.

\begin{table}[H]
  \centering
  \caption{Parámetros de \texttt{DynamicsCPU.jac}}
  \small
  \begin{tabularx}{\textwidth}{@{}p{0.27\textwidth}p{0.33\textwidth}X@{}}
    \toprule
    \textbf{Parámetro} & \textbf{Tipo} & \textbf{Descripción} \\
    \midrule
    \texttt{x} & \texttt{array-like (8,)} & Estado en el que se evalúa el Jacobiano. \\
    \texttt{t} & \texttt{float} & Tiempo de evaluación. \\
    \texttt{m1, m2} & \texttt{float} & Masas de los cuerpos. \\
    \texttt{G} & \texttt{float} & Constante gravitacional. \\
    \bottomrule
  \end{tabularx}
\end{table}

\begin{table}[H]
  \centering
  \caption{Valor de retorno de \texttt{DynamicsCPU.jac}}
  \small
  \begin{tabularx}{\textwidth}{@{}p{0.35\textwidth}X@{}}
    \toprule
    \textbf{Tipo} & \textbf{Descripción} \\
    \midrule
    \texttt{numpy.ndarray} & Matriz \(8\times8\) con derivadas parciales \(\partial f/\partial x\). \\
    \bottomrule
  \end{tabularx}
\end{table}

\begin{listing}[H]
\caption{Bloque de Código de \texttt{DynamicsCPU.jac}}
\begin{minted}[fontsize=\small, breaklines, breakanywhere]{python}
class DynamicsCPU:
    """Calcula el RHS y jacobiano del sistema de dos cuerpos en 2D."""
        def jac(self, x: Any, t: float, m1: float, m2: float, G: float) -> Any:
        import numpy as _np

        eps = 1e-12
        x1, y1, vx1, vy1, x2, y2, vx2, vy2 = x
        dx = x2 - x1
        dy = y2 - y1
        r2 = dx * dx + dy * dy + eps
        r = _np.sqrt(r2)
        r3 = r2 * r
        r5 = r3 * r2

        dfx_dx1 = G * (m2 * (3.0 * dx * dx / r5 - 1.0 / r3))
        dfx_dy1 = G * (m2 * (3.0 * dx * dy / r5))
        dfx_dx2 = -dfx_dx1
        dfx_dy2 = -dfx_dy1

        dfy_dx1 = G * (m2 * (3.0 * dx * dy / r5))
        dfy_dy1 = G * (m2 * (3.0 * dy * dy / r5 - 1.0 / r3))
        dfy_dx2 = -dfy_dx1
        dfy_dy2 = -dfy_dy1

        g2x_dx1 = G * (m1 * (3.0 * dx * dx / r5 - 1.0 / r3))
        g2x_dy1 = G * (m1 * (3.0 * dx * dy / r5))
        g2x_dx2 = -g2x_dx1
        g2x_dy2 = -g2x_dy1

        g2y_dx1 = G * (m1 * (3.0 * dx * dy / r5))
        g2y_dy1 = G * (m1 * (3.0 * dy * dy / r5 - 1.0 / r3))
        g2y_dx2 = -g2y_dx1
        g2y_dy2 = -g2y_dy1

        J = _np.zeros((8, 8), dtype=float)
        J[0, 2] = 1.0
        J[1, 3] = 1.0
        J[4, 6] = 1.0
        J[5, 7] = 1.0

        J[2, 0] = -dfx_dx1
        J[2, 1] = -dfx_dy1
        J[2, 4] = dfx_dx2
        J[2, 5] = dfx_dy2

        J[3, 0] = -dfy_dx1
        J[3, 1] = -dfy_dy1
        J[3, 4] = dfy_dx2
        J[3, 5] = dfy_dy2

        J[6, 0] = g2x_dx1
        J[6, 1] = g2x_dy1
        J[6, 4] = g2x_dx2
        J[6, 5] = g2x_dy2

        J[7, 0] = g2y_dx1
        J[7, 1] = g2y_dy1
        J[7, 4] = g2y_dx2
        J[7, 5] = g2y_dy2

        return J
\end{minted}
\end{listing}

\begin{listing}[H]
\caption{Ejemplo de uso de \texttt{DynamicsCPU.jac}}
\begin{minted}[fontsize=\small, breaklines, breakanywhere]{python}
J = dyn.jac(state, t=0.0, m1=1.0, m2=1.0, G=1.0)  # shape: (8, 8)
\end{minted}
\end{listing}

\subsubsection{\texttt{class DynamicsGPU}}
\paragraph{Propósito}
Implementación acelerada en GPU del RHS \( \dot{x}=f(x,t) \) y del Jacobiano \(J=\partial f/\partial x\) para el problema de dos cuerpos en 2D usando \textbf{CuPy}. Detecta disponibilidad en \texttt{\_\_init\_\_}, y exige GPU vía \texttt{\_require\_gpu~()}. El estado sigue el mismo orden que en CPU:\
\[
x = (x_1, y_1, v_{x1}, v_{y1},\; x_2, y_2, v_{x2}, v_{y2}).
\]
Se usa un suavizado numérico \(\varepsilon = 10^{-12}\) para evitar divisiones por cero en \(r^{-3}\) y \(r^{-5}\).

\begin{table}[H]
  \centering
  \caption{\texttt{DynamicsGPU}: interfaz}
  \small
  \begin{tabularx}{\textwidth}{@{}p{0.27\textwidth}p{0.33\textwidth}X@{}}
    \toprule
    \textbf{Miembro} & \textbf{Tipo/Firma} & \textbf{Descripción} \\
    \midrule
    \texttt{available} & \texttt{bool} & \texttt{True} si \texttt{cupy} está disponible en el entorno. \\
    \texttt{\_require\_gpu} & \texttt{\_require\_gpu~() -> None} & Lanza \texttt{RuntimeError} si no hay GPU/CuPy. \\
    \texttt{f} & \texttt{f~(x, t, m1, m2, G) -> cp.ndarray} & Campo vectorial (dim. 8). \\
    \texttt{jac} & \texttt{jac~(x, t, m1, m2, G) -> cp.ndarray} & Jacobiano (matriz \(8\times 8\)). \\
    \bottomrule
  \end{tabularx}
\end{table}

\subsubsection*{\texttt{f~(x, t, m1, m2, G) -> cp.ndarray}}
\textbf{Propósito.} Evaluar el RHS en GPU para dos cuerpos en 2D, devolviendo \((\dot{x}_1,\dot{y}_1,\dot{v}_{x1},\dot{v}_{y1},\dot{x}_2,\dot{y}_2,\dot{v}_{x2},\dot{v}_{y2})\).

\begin{table}[H]
  \centering
  \caption{Parámetros de \texttt{DynamicsGPU.f}}
  \small
  \begin{tabularx}{\textwidth}{@{}p{0.27\textwidth}p{0.33\textwidth}X@{}}
    \toprule
    \textbf{Parámetro} & \textbf{Tipo} & \textbf{Descripción} \\
    \midrule
    \texttt{x} & \texttt{array-like (8,)} & Estado: \((x_1,y_1,v_{x1},v_{y1},x_2,y_2,v_{x2},v_{y2})\). \\
    \texttt{t} & \texttt{float} & Tiempo de evaluación (no autónomo opcional). \\
    \texttt{m1, m2} & \texttt{float} & Masas de los cuerpos. \\
    \texttt{G} & \texttt{float} & Constante gravitacional. \\
    \bottomrule
  \end{tabularx}
\end{table}

\begin{table}[H]
  \centering
  \caption{Valor de retorno de \texttt{DynamicsGPU.f}}
  \small
  \begin{tabularx}{\textwidth}{@{}p{0.35\textwidth}X@{}}
    \toprule
    \textbf{Tipo} & \textbf{Descripción} \\
    \midrule
    \texttt{cupy.ndarray (8,)} & Derivadas del estado, \texttt{float64}. \\
    \bottomrule
  \end{tabularx}
\end{table}

\begin{listing}[H]
\caption{Bloque de código de \texttt{DynamicsGPU.f}}
\begin{minted}[fontsize=\small, breaklines, breakanywhere]{python}
class DynamicsGPU:
    """Implementación del RHS y jacobiano en GPU usando CuPy."""

    def __init__(self) -> None:
        self.available = False
        try:
            import cupy as _cp  # noqa: F401
            self.available = True
        except Exception:
            self.available = False

    def _require_gpu(self) -> None:
        if not self.available:
            raise RuntimeError(
                "CuPy no está disponible. Instálalo o usa DynamicsCPU."
            )

    def f(self, x: Any, t: float, m1: float, m2: float, G: float) -> Any:
        self._require_gpu()
        import cupy as cp

        eps = cp.float64(1e-12)
        state = cp.asarray(x, dtype=cp.float64)

        x1, y1, vx1, vy1, x2, y2, vx2, vy2 = state
        dx = x2 - x1
        dy = y2 - y1
        r2 = dx * dx + dy * dy + eps
        r = cp.sqrt(r2)
        r3 = r2 * r

        ax1 = G * m2 * dx / r3
        ay1 = G * m2 * dy / r3
        ax2 = -G * m1 * dx / r3
        ay2 = -G * m1 * dy / r3

        return cp.array([vx1, vy1, ax1, ay1, vx2, vy2, ax2, ay2], dtype=cp.float64)
\end{minted}
\end{listing}

\begin{listing}[H]
\caption{Ejemplo de uso de \texttt{DynamicsGPU.f} con retroceso a CPU}
\begin{minted}[fontsize=\small, breaklines, breakanywhere]{python}
import numpy as np
from simulacion.dynamics import DynamicsGPU, DynamicsCPU

state = np.array([-1.0, 0.0, 0.0, -0.5, 1.0, 0.0, 0.0, 0.5], dtype=float)
m1 = m2 = G = 1.0
t = 0.0

dyn_gpu = DynamicsGPU()
try:
    rhs_gpu = dyn_gpu.f(state, t=t, m1=m1, m2=m2, G=G)  # cupy.ndarray
    print(rhs_gpu.shape)
except RuntimeError:
    # Fallback si no hay CuPy
    dyn_cpu = DynamicsCPU()
    rhs_cpu = dyn_cpu.f(state, t=t, m1=m1, m2=m2, G=G)
    print(rhs_cpu.shape)
\end{minted}
\end{listing}

\subsubsection*{\texttt{jac~(x, t, m1, m2, G) -> cp.ndarray}}
\textbf{Propósito.} Calcular en GPU la matriz Jacobiana \(8\times 8\), incluyendo derivadas cruzadas \((x_1,y_1)\)–\((x_2,y_2)\) según las expresiones analíticas de \(\partial (a_{x}, a_{y})/\partial (x,y)\) para potencial newtoniano en 2D.

\begin{table}[H]
  \centering
  \caption{Parámetros de \texttt{DynamicsGPU.jac}}
  \small
  \begin{tabularx}{\textwidth}{@{}p{0.27\textwidth}p{0.33\textwidth}X@{}}
    \toprule
    \textbf{Parámetro} & \textbf{Tipo} & \textbf{Descripción} \\
    \midrule
    \texttt{x} & \texttt{array-like (8,)} & Estado donde se evalúa el Jacobiano. \\
    \texttt{t} & \texttt{float} & Tiempo de evaluación. \\
    \texttt{m1, m2} & \texttt{float} & Masas de los cuerpos. \\
    \texttt{G} & \texttt{float} & Constante gravitacional. \\
    \bottomrule
  \end{tabularx}
\end{table}

\begin{table}[H]
  \centering
  \caption{Valor de retorno de \texttt{DynamicsGPU.jac}}
  \small
  \begin{tabularx}{\textwidth}{@{}p{0.35\textwidth}X@{}}
    \toprule
    \textbf{Tipo} & \textbf{Descripción} \\
    \midrule
    \texttt{cupy.ndarray (8, 8)} & Jacobiano del sistema, \texttt{float64}. \\
    \bottomrule
  \end{tabularx}
\end{table}

\begin{listing}[H]
\caption{Bloque de Código de \texttt{DynamicsGPU.jac}}
\begin{minted}[fontsize=\small, breaklines, breakanywhere]{python}
class DynamicsGPU:
    """Implementación del RHS y jacobiano en GPU usando CuPy."""
    def jac(self, x: Any, t: float, m1: float, m2: float, G: float) -> Any:
        self._require_gpu()
        import cupy as cp

        eps = cp.float64(1e-12)
        state = cp.asarray(x, dtype=cp.float64)

        x1, y1, vx1, vy1, x2, y2, vx2, vy2 = state
        dx = x2 - x1
        dy = y2 - y1
        r2 = dx * dx + dy * dy + eps
        r = cp.sqrt(r2)
        r3 = r2 * r
        r5 = r3 * r2

        dfx_dx1 = G * (m2 * (3.0 * dx * dx / r5 - 1.0 / r3))
        dfx_dy1 = G * (m2 * (3.0 * dx * dy / r5))
        dfx_dx2 = -dfx_dx1
        dfx_dy2 = -dfx_dy1

        dfy_dx1 = G * (m2 * (3.0 * dx * dy / r5))
        dfy_dy1 = G * (m2 * (3.0 * dy * dy / r5 - 1.0 / r3))
        dfy_dx2 = -dfy_dx1
        dfy_dy2 = -dfy_dy1

        g2x_dx1 = G * (m1 * (3.0 * dx * dx / r5 - 1.0 / r3))
        g2x_dy1 = G * (m1 * (3.0 * dx * dy / r5))
        g2x_dx2 = -g2x_dx1
        g2x_dy2 = -g2x_dy1

        g2y_dx1 = G * (m1 * (3.0 * dx * dy / r5))
        g2y_dy1 = G * (m1 * (3.0 * dy * dy / r5 - 1.0 / r3))
        g2y_dx2 = -g2y_dx1
        g2y_dy2 = -g2y_dy1

        J = cp.zeros((8, 8), dtype=cp.float64)
        J[0, 2] = 1.0
        J[1, 3] = 1.0
        J[4, 6] = 1.0
        J[5, 7] = 1.0

        J[2, 0] = -dfx_dx1
        J[2, 1] = -dfx_dy1
        J[2, 4] = dfx_dx2
        J[2, 5] = dfx_dy2

        J[3, 0] = -dfy_dx1
        J[3, 1] = -dfy_dy1
        J[3, 4] = dfy_dx2
        J[3, 5] = dfy_dy2

        J[6, 0] = g2x_dx1
        J[6, 1] = g2x_dy1
        J[6, 4] = g2x_dx2
        J[6, 5] = g2x_dy2

        J[7, 0] = g2y_dx1
        J[7, 1] = g2y_dy1
        J[7, 4] = g2y_dx2
        J[7, 5] = g2y_dy2

        return J
\end{minted}
\end{listing}

\begin{listing}[H]
\caption{Ejemplo de uso de \texttt{DynamicsGPU.jac}}
\begin{minted}[fontsize=\small, breaklines, breakanywhere]{python}
import numpy as np
from simulacion.dynamics import DynamicsGPU

x = np.array([-1.0, 0.0, 0.0, -0.5, 1.0, 0.0, 0.0, 0.5], dtype=float)
J = DynamicsGPU().jac(x, t=0.0, m1=1.0, m2=1.0, G=1.0)  # cupy.ndarray (8, 8)
\end{minted}
\end{listing}

% ---------- lyapunov.py ----------
\subsection{\texttt{lyapunov.py} — Estimación de Exponente de Lyapunov}
\paragraph{Descripción general}
\texttt{LyapunovEstimator} expone dos rutas de cálculo del mLCE:\ una corta (\texttt{mLCE\_short}) y una larga (\texttt{mLCE\_long}), ambas soportadas por un método interno que intenta, primero, la API variacional de REBOUND y, en su defecto, un esquema discreto de Benettin con re-normalización.

\subsubsection{\texttt{class LyapunovEstimator}}
\paragraph{Propósito} Estimar el exponente de Lyapunov máximo (mLCE) a partir de un contexto de simulación REBOUND o de una trayectoria discreta.

\subsubsection*{\texttt{mLCE\_short~(trajectory=None, x0=None, window=100.0) -> dict}}
\textbf{Propósito.} Estimar mLCE sobre una ventana temporal corta.

\begin{table}[H]
  \centering
  \caption{Parámetros de \texttt{mLCE\_short}}
  \small
  \begin{tabularx}{\textwidth}{@{}p{0.27\textwidth}p{0.33\textwidth}X@{}}
    \toprule
    \textbf{Parámetro} & \textbf{Tipo} & \textbf{Descripción} \\
    \midrule
    \texttt{trajectory} & \texttt{Any} & Simulación REBOUND o contexto discreto. \\
    \texttt{x0} & \texttt{Iterable[float] | None} & (Reservado) Semilla/estado inicial, si aplica. \\
    \texttt{window} & \texttt{float} & Horizonte temporal (unidades de la simulación). \\
    \bottomrule
  \end{tabularx}
\end{table}

\begin{table}[H]
  \centering
  \caption{Valor de retorno de \texttt{mLCE\_short}}
  \small
  \begin{tabularx}{\textwidth}{@{}p{0.35\textwidth}X@{}}
    \toprule
    \textbf{Tipo} & \textbf{Descripción} \\
    \midrule
    \texttt{dict} & \{\texttt{``lambda''}, \texttt{``series''}, \texttt{``meta''}\} con detalles de cómputo. \\
    \bottomrule
  \end{tabularx}
\end{table}

\begin{listing}[H]
\caption{Bloque de Código de \texttt{mLCE\_short}}
\begin{minted}[fontsize=\small, breaklines, breakanywhere]{python}
class LyapunovEstimator:
    def __init__(self) -> None:
        self._cpu_dyn = DynamicsCPU()

    def mLCE_short(
        self,
        trajectory: Any = None,
        x0: Optional[Iterable[float]] = None,
        window: float = 100.0,
    ) -> Dict[str, Any]:
        try:
            lam, series, meta = self._mlce_with_context(trajectory, window)
            return {"lambda": lam, "series": series, "meta": meta}
        except Exception as e:
            return {"lambda": float("inf"), "series": None, "meta": {"impl": "error", "err": str(e)}}
\end{minted}
\end{listing}

\begin{listing}[H]
\caption{Ejemplo de uso de \texttt{mLCE\_short}}
\begin{minted}[fontsize=\small, breaklines, breakanywhere]{python}
from simulacion.rebound_adapter import ReboundSim
from simulacion.lyapunov import LyapunovEstimator

sim_api = ReboundSim(G=1.0, integrator="whfast")
sim = sim_api.setup_simulation(
    masses=(1.0, 1.0),
    r0=(( -1.0, 0.0, 0.0), (1.0, 0.0, 0.0)),
    v0=((  0.0,-0.5, 0.0), (0.0, 0.5, 0.0)),
)
est = LyapunovEstimator()
res = est.mLCE_short(trajectory=sim, window=100.0)
print(res["lambda"], res["meta"])
\end{minted}
\end{listing}

\subsubsection*{\texttt{mLCE\_long~(trajectory=None, x0=None, window=1000.0) -> dict}}
\textbf{Propósito.} Estimar mLCE sobre una ventana temporal mayor (muestras más estables para diagnóstico).

\begin{table}[H]
  \centering
  \caption{Parámetros y retorno de \texttt{mLCE\_long}}
  \small
  \begin{tabularx}{\textwidth}{@{}p{0.27\textwidth}p{0.33\textwidth}X@{}}
    \toprule
    \textbf{Parámetro} & \textbf{Tipo} & \textbf{Descripción} \\
    \midrule
    \texttt{trajectory} & \texttt{Any} & Simulación REBOUND o contexto discreto. \\
    \texttt{x0} & \texttt{Iterable[float] | None} & (Reservado) Semilla/estado inicial, si aplica. \\
    \texttt{window} & \texttt{float} & Horizonte temporal (unidades de la simulación). \\
    \bottomrule
  \end{tabularx}

  \vspace{0.5em}
  \begin{tabularx}{\textwidth}{@{}p{0.35\textwidth}X@{}}
    \toprule
    \textbf{Tipo de retorno} & \textbf{Descripción} \\
    \midrule
    \texttt{dict} & \{\texttt{``lambda''}, \texttt{``series''}, \texttt{``meta''}\}. \\
    \bottomrule
  \end{tabularx}
\end{table}

\begin{listing}[H]
\caption{Bloque de Código de \texttt{mLCE\_long}}
\begin{minted}[fontsize=\small, breaklines, breakanywhere]{python}
class LyapunovEstimator:

    def mLCE_long(
        self,
        trajectory: Any = None,
        x0: Optional[Iterable[float]] = None,
        window: float = 1000.0,
    ) -> Dict[str, Any]:
        try:
            lam, series, meta = self._mlce_with_context(trajectory, window)
            return {"lambda": lam, "series": series, "meta": meta}
        except Exception as e:
            return {"lambda": float("inf"), "series": None, "meta": {"impl": "error", "err": str(e)}}
\end{minted}
\end{listing}

\begin{listing}[H]
\caption{Ejemplo de uso de \texttt{mLCE\_long}}
\begin{minted}[fontsize=\small, breaklines, breakanywhere]{python}
res_long = est.mLCE_long(trajectory=sim, window=1000.0)
print(res_long["lambda"])
\end{minted}
\end{listing}

\subsubsection*{Notas sobre la implementación interna}
Cuando la API variacional de REBOUND está disponible, se añade al simulador un sistema variacional y se integra hasta completar el horizonte \texttt{window}, extrayendo el LE.\ Si la API no existe o falla, se replica el sistema para un par \emph{referencia/perturbado}, se aplica una perturbación pequeña \(\varepsilon\) al estado y, tras cada paso, se re-normaliza la separación para acumular el promedio logarítmico \(\lambda\).

\begin{listing}[H]
\caption{Bloque de Código de métodos auxiliares I}
\begin{minted}[fontsize=\small, breaklines, breakanywhere]{python}
class LyapunovEstimator:

    def _mlce_with_context(self, ctx: Any, window: float) -> Tuple[float, Optional[Any], Dict[str, Any]]:
        import numpy as _np

        try:
            import rebound  # noqa: F401
        except Exception:
            raise

        if ctx is None:
            raise ValueError("Se requiere una simulacion o contexto discreto para estimar Lyapunov")

        if isinstance(ctx, dict) and "sim" in ctx:
            sim = ctx["sim"]
            dt = float(ctx.get("dt", 0.5))
            t_end = float(ctx.get("t_end", window))
            masses_ctx = ctx.get("masses")
        else:
            sim = ctx
            dt = 0.5
            t_end = window
            masses_ctx = None

        if masses_ctx is not None:
            masses = tuple(float(m) for m in masses_ctx)
        else:
            masses, _, _ = self._extract_rebound_state(sim)

        if dt <= 0.0:
            raise ValueError("dt debe ser positivo para el Lyapunov discreto")

        steps = max(1, math.ceil(t_end / dt)) if t_end is not None else 1

        try:
            lam = self._mlce_rebound_variational(sim, dt=dt, steps=steps)
            meta = {"impl": "rebound_variational_discrete", "steps": steps, "dt": dt, "n_bodies": len(masses), "masses": masses}
            return lam, None, meta
        except Exception:
            lam, series = self._mlce_benettin(sim, masses, dt=dt, steps=steps)
            meta = {"impl": "benettin_discrete", "steps": steps, "dt": dt, "n_bodies": len(masses), "masses": masses}
            return lam, series, meta
\end{minted}
\end{listing}

\begin{listing}[H]
\caption{Bloque de Código de métodos auxiliares II}
\begin{minted}[fontsize=\small, breaklines, breakanywhere]{python}
class LyapunovEstimator:

    @staticmethod
    def _extract_rebound_state(
        sim: Any,
    ) -> Tuple[Tuple[float, ...], Tuple[Tuple[float, float, float], ...], Tuple[Tuple[float, float, float], ...]]:
        if not hasattr(sim, "particles"):
            raise ValueError("La simulacion proporcionada no contiene particulas.")
        masses: list[float] = []
        positions: list[Tuple[float, float, float]] = []
        velocities: list[Tuple[float, float, float]] = []
        for p in sim.particles:
            masses.append(float(getattr(p, "m", 0.0)))
            positions.append((float(p.x), float(p.y), float(p.z)))
            velocities.append((float(p.vx), float(p.vy), float(p.vz)))
        return tuple(masses), tuple(positions), tuple(velocities)
\end{minted}
\end{listing}

\begin{listing}[H]
\caption{Bloque de Código de métodos auxiliares III}
\begin{minted}[fontsize=\small, breaklines, breakanywhere]{python}
class LyapunovEstimator:

    def _mlce_rebound_variational(self, sim: Any, dt: float, steps: int) -> float:
        if not hasattr(sim, "particles"):
            raise RuntimeError("Invalid REBOUND Simulation provided")
        try:
            if hasattr(sim, "add_variationals"):
                sim.add_variationals(1)
            elif hasattr(sim, "add_variational"):
                sim.add_variational()
            else:
                raise AttributeError("REBOUND variational API not found")

            sim.dt = dt
            steps = max(1, int(steps))
            if hasattr(sim, "calculate_lyapunov"):
                t = sim.t if hasattr(sim, "t") else 0.0
                sim.integrate(t + steps * dt)
                val = sim.calculate_lyapunov()
                return float(val) * dt
            lam_sum = 0.0
            d0 = None
            base_t = sim.t if hasattr(sim, "t") else 0.0
            for i in range(steps):
                sim.integrate(base_t + (i + 1) * dt)
                if hasattr(sim, "variationals") and len(sim.variationals) > 0:
                    try:
                        import numpy as _np

                        vv = _np.asarray(list(sim.variationals[0]))
                        d = float(_np.linalg.norm(vv))
                    except Exception:
                        raise AttributeError("Unsupported variational vector structure")
                else:
                    raise AttributeError("variationals not accessible")
                if d0 is None:
                    d0 = max(d, 1e-30)
                lam_sum += _np.log(max(d, 1e-30) / d0)
            return lam_sum / float(steps)
        except Exception as e:
            raise RuntimeError(f"REBOUND variational path failed: {e}")
\end{minted}
\end{listing}

\begin{listing}[H]
\caption{Bloque de Código de métodos auxiliares IV}
\begin{minted}[fontsize=\small, breaklines, breakanywhere]{python}
class LyapunovEstimator:

    def _mlce_benettin(self, sim: Any, masses: Tuple[float, ...], dt: float, steps: int) -> Tuple[float, Any]:
        import numpy as _np

        masses = tuple(float(m) for m in masses)
        if len(masses) < 2:
            raise ValueError("Se requieren al menos dos masas para estimar el Lyapunov discreto.")

        try:
            sim_ref = sim.copy()
            sim_pert = sim.copy()
        except Exception:
            extracted = self._extract_rebound_state(sim)
            masses_extracted, r0_init, v0_init = extracted
            rbs = ReboundSim(G=getattr(sim, "G", 1.0), integrator=getattr(sim, "integrator", "whfast"))
            sim_ref = rbs.setup_simulation(masses_extracted, r0_init, v0_init, move_to_com=False)
            sim_pert = rbs.setup_simulation(masses_extracted, r0_init, v0_init, move_to_com=False)

        eps = 1e-9
        sim_ref.dt = dt
        sim_pert.dt = dt
        sim_pert.particles[0].x += eps

        def state_vector(S: Any) -> _np.ndarray:
            data = []
            for p in S.particles:
                data.extend([p.x, p.y, p.z, p.vx, p.vy, p.vz])
            return _np.array(data, dtype=float)

        def set_state(S: Any, vec: _np.ndarray) -> None:
            for idx, p in enumerate(S.particles):
                base = idx * 6
                p.x = float(vec[base])
                p.y = float(vec[base + 1])
                p.z = float(vec[base + 2])
                p.vx = float(vec[base + 3])
                p.vy = float(vec[base + 4])
                p.vz = float(vec[base + 5])

        x_ref0 = state_vector(sim_ref)
        x_pert0 = state_vector(sim_pert)
        d0 = float(_np.linalg.norm(x_pert0 - x_ref0))
        d0 = max(d0, 1e-30)

        steps = max(1, int(steps))
        lam_sum = 0.0
        series = _np.zeros(steps, dtype=float)

        base_t = getattr(sim_ref, "t", 0.0)

        for i in range(steps):
            t_target = base_t + (i + 1) * dt
            sim_ref.integrate(t_target)
            sim_pert.integrate(t_target)

            xr = state_vector(sim_ref)
            xp = state_vector(sim_pert)
            diff = xp - xr
            d = float(_np.linalg.norm(diff))
            d = max(d, 1e-30)
            lam_sum += _np.log(d / d0)
            series[i] = lam_sum / float(i + 1)

            alpha = d0 / d
            new_xp = xr + alpha * diff
            set_state(sim_pert, new_xp)

        lam = lam_sum / float(steps)
        return lam, series
            return {"lambda": float("inf"), "series": None, "meta": {"impl": "error", "err": str(e)}}
\end{minted}
\end{listing}
% ---------- __init__.py ----------
\subsection{\texttt{\_\_init\_\_.py} — API del Paquete}
\paragraph{Descripción general}
Expone una interfaz pública compacta para importar las clases principales del módulo:
\begin{itemize}
    \item \texttt{DynamicsCPU}, \texttt{DynamicsGPU}
    \item \texttt{ReboundSim}
    \item \texttt{LyapunovEstimator}
\end{itemize}

\begin{listing}[H]
\caption{Bloque de Código de \texttt{mLCE\_long}}
\begin{minted}[fontsize=\small, breaklines, breakanywhere]{python}
"""
Capa de simulacion que encapsula REBOUND y el calculo del exponente de Lyapunov.
"""

from .dynamics import DynamicsCPU, DynamicsGPU  # noqa: F401
from .rebound_adapter import ReboundSim  # noqa: F401
from .lyapunov import LyapunovEstimator  # noqa: F401

__all__ = ["DynamicsCPU", "DynamicsGPU", "ReboundSim", "LyapunovEstimator"]


if __name__ == "__main__":
    dyn = DynamicsCPU()
    print("Norma del Jacobiano con masas=1:", dyn.jac([1, 0, 0, 0, -1, 0, 0, 0], 0, 1.0, 1.0, 1.0))

\end{minted}
\end{listing}

\begin{listing}[H]
\caption{Ejemplo de importación desde el paquete}
\begin{minted}[fontsize=\small, breaklines, breakanywhere]{python}
from simulacion import DynamicsCPU, ReboundSim, LyapunovEstimator
\end{minted}
\end{listing}
% --- Fin de capítulo ---
